\paragraph{Prescribed finite phase coherence.}
A fluctuating drive phase can be specified without identifying it with
thermal noise. Let the finite-duration source be
$S_\theta(t)=e^{i\theta(t)}S_d(t)$, where the initial phase is uniform and
independent and $\theta(t)-\theta(\Delta)=\sqrt{2D}\,B_{t-\Delta}$ during
the pulse, with $B$ a standard real Wiener process independent of that initial phase. Then
$\mathbb E e^{i[\theta(t)-\theta(s)]}=e^{-D|t-s|}$ and the anomalous
source correlation vanishes. Every realization retains the source's
$L^2$ budget. This does not fix the work delivered to the field.

For the fixed finite real-linear response used above, let $m(\delta)$ denote
the phase-averaged source work when the \emph{entire} prescribed source
is multiplied by $e^{i\delta(t-\Delta)}$. Its phase-diffusion mean obeys
\[
 m_D=\int_{-\infty}^{\infty}
 \frac{D}{\pi(D^2+\delta^2)}\,m(\delta)\,d\delta,
 \qquad D>0.
\]
This follows from the characteristic function of the Cauchy distribution
and the quadratic response kernel. At finite time the finite-dimensional
propagator is bounded; compact source support and finite $L^2$ norm give
absolute integrability and justify interchanging the integrals.
The same weighting convolves the finite-pulse spectral power and preserves
its integrated weight. Since the weight is even, the odd detuning response
cancels. The positive local slope found above therefore does not determine
whether phase diffusion increases or decreases mean work. Nor does a single
symmetric detuning pair settle this question.

For computation of this mean work, the lag factor $e^{-D(t-s)}$ can instead
be included in the auxiliary generator $A-DI$, evaluated for two genuine
source quadratures. With state coordinates $(q,p)$, both equations acquire
$-D$ terms; the work output remains $p+i\omega q$, although now
$\dot q=p-Dq$. These auxiliary states are not stochastic field trajectories,
and their endpoint energies are not ensemble energies. This reduction
concerns the specified work observable and does not insert physical friction
into the original field equations.

A Cauchy-distributed constant detuning and a Wiener phase process share the
relevant two-point statistics but not their trajectory laws. Their quadratic
response means agree; the auxiliary damped endpoint is a different object.
Ideal Wiener phase also has no finite ultraviolet cutoff. Thus this identity
is a finite-model interpretation, not a continuum certificate, a temperature
assignment, or a derivation of an environmental bath. Physical bandwidth and
coupling require separate specification before experimental use.
\paragraph{One finite-coherence comparison.}
We fixed $D=1$, corresponding to one pulse duration for the exponential
phase-correlation time, and compared it with $D=0$. The two genuine
quadratures of each auxiliary response were evolved on the same two
discretizations as the detuning test. The finer grid gives mean source
works $m_0=1.133373236\times10^{-4}$ and
$m_1=2.173314375\times10^{-4}$. Their difference,
$1.039941139\times10^{-4}$, is positive on both grids and exceeds the
predeclared diagnostic guard $6.63\times10^{-9}$. The undiffused adapter
reproduces the earlier work values, and the modified auxiliary balances
include $-2D$ times each quadratic form. Free and restoring-oscillator
controls test both the work output and the placement of damping in both
auxiliary coordinates.

This is an increase in mean work supplied by the prescribed external
source at fixed source $L^2$ budget. It is not a gain in energetic efficiency
or a demonstration of self-sustained oscillations. The guard combines
grid differences and implementation residuals; it is not a certified
error bound, and space and time resolution change together. The nominal
profile, finite boundary and linear-response limitations persist. Only
one nonzero diffusion constant is compared: no noise optimum, temperature,
nonlinear formation or lifetime conclusion is established.
The same source-output reduction with output $p$ gives an additional
observable. In the undamped physical finite model, the rotating quadratic
generator satisfies $\dot H=2\operatorname{Re}\langle p,S_\theta\rangle$.
For zero initial perturbation, its ensemble endpoint mean therefore equals
the half-sum of the two auxiliary integrals of this \emph{original} source
rate. It does not equal the auxiliary endpoint $H$. Extracting these
already stored integrals gives, on the finer grid,
$\mathbb E H(T)=2.09020168\times10^{-4}$ for $D=0$ and
$2.93069596\times10^{-4}$ for $D=1$. Both grids show this approximately
$40\%$ increase; this is a descriptive follow-up without a new certified
comparison bound or field integration.

This rotating quadratic generator need not be positive, and its global
value does not measure energy localized near the Q-ball. Full laboratory
endpoint energy additionally requires the physical ensemble charge,
including the anomalous state-exchange contribution. The latter is not
given by the exchange evaluated on the damped auxiliary states. Thus
neither the source-work increase nor this rotating-generator mean establishes
laboratory energy retention or improved confinement.
